\section{Chronology-induced real multiplication}\label{sec:rm}

Let $q:E\to C$ be the reflection quotient.  The two maps
\[
 E\overset{q}{\longrightarrow}C,
 \qquad
 E\overset{q\circ\tau^{-1}}{\longrightarrow}C
\]
define a degree-$(2,2)$ correspondence on $C$.  Its unnormalized push-pull on
the Jacobian is
\begin{equation}\label{eq:A}
 A=q_*\tau^*q^*.
\end{equation}
The word ``unnormalized'' is important: the idempotent
$e_J=(1+J)/2$ gives $e_J\tau e_J=A/2$ on scalar cohomology and therefore a
rescaled polynomial.

\begin{theorem}[H\'enon real multiplication]\label{thm:rm}
The endomorphism $A$ is Rosati self-adjoint and has exact minimal polynomial
\begin{equation}\label{eq:minpoly}
 m(X)=X^3+X^2-2X-1.
\end{equation}
Consequently
\[
 \boxed{
 \Q(\zeta_7+\zeta_7^{-1})
 \hookrightarrow \End_\Q^0(\Jac(C)).
 }
\]
\end{theorem}

\begin{proof}
Put $V=H^1(E,\Q)$ and $H=\langle J\rangle$.  Pullback identifies
$H^1(C,\Q)$ with $V^H$, and $q^*q_*=1+J$ on $V$.  For $v\in V^H$,
\[
 q^*A(q^*)^{-1}v
 =(1+J)\tau v
 =\tau v+\tau^{-1}v.
\]
Thus $A$ is the adjacency operator $X=\tau+\tau^{-1}$ on the reflection-fixed
subspace.

There is no rotation-fixed vector in $V^H$: such a vector would be fixed by
all of $D_7$, whereas
$V^{D_7}=H^1(E/D_7,\Q)=H^1(\Pone,\Q)=0$.  Hence
$1+\tau+\cdots+\tau^6=0$ on every relevant eigenspace.  Pair inverse powers
and use
\[
 \tau^2+\tau^{-2}=X^2-2,
 \qquad
 \tau^3+\tau^{-3}=X^3-3X
\]
to obtain $m(X)=0$.

It remains to prove that the cubic is minimal, not merely annihilating.  Let
$\varepsilon$ be the reflection sign representation and let $W$ be the
six-dimensional rational irreducible whose complexification is the sum of
the three nontrivial two-dimensional representations.  Quotient genera give
\[
 H^1(E,\Q)\cong\varepsilon^{\oplus4}\oplus W^{\oplus2}.
\]
Indeed, there are no trivial summands; $C_7$-invariants have dimension
$2g(B)=4$; and reflection invariants have dimension $2g(C)=6$.  On $W^H$,
$X$ has all three distinct eigenvalues
\[
 \zeta_7^k+\zeta_7^{-k},\qquad k=1,2,3,
\]
each twice on $H^1(C)$.  Thus the characteristic polynomial of $A$ is
$m^2$ and its minimal polynomial is the irreducible cubic $m$, whose
discriminant is 49.

Finally, the Rosati adjoint of Eq.~\eqref{eq:A} is
$q_*(\tau^{-1})^*q^*$.  Since $\tau^{-1}=J\tau J$ and $q\circ J=q$, this is
again $A$.  The generated field is totally real and fixed pointwise by
Rosati, proving the stated real multiplication.
\end{proof}

This is precisely the double-coset mechanism of
Ellenberg~\cite{ellenberg2001endomorphism}; genus-three realizations with the
same cubic were constructed explicitly in~\cite{hoffman2016genus3}.  The
theorem identifies the H\'enon time step as the generator for the present
curve.  We do not prove that this cubic field is the entire endomorphism
algebra or that $\Jac(C)$ is absolutely simple.

The complex $\tau$-character dimensions follow at once:
\[
 \dim H^1(E)_{k=0}=4,
 \qquad
 \dim H^1(E)_{k}=2\quad(1\le k\le6),
\]
and reversal pairs $k$ with $-k$.  These dimensions are group-forced; the
arithmetic content lies in the rank-two Frobenius action inside each paired
sector.
